%% kapa: exam=1 round=19 year=2008 period=1 form=A subject=civil_law subject_ko=민법 seq=12 number=12
%% source: Q-Net 감정평가사 기출문제 (cst00309) artlSeq=1001962 file=extracted/1008/1001962/2008년 제19회 감정평가사 1차시험 기출문제/1교시.hwp
%% conversion: kordoc 4.24.51 -> tools/gen1.py (auto); review=proofread-sonnet; flags=none
\begin{question}{12}
소멸시효의 기산점에 관한 설명으로 틀린 것은? (다툼이 있는 경우에는 판례에 의함)
\begin{choices}
  \item 권리자의 개인적 사정이나 법률지식의 부족, 권리존재의 부지 또는 채무자의 부재 등 사실상 장애는 소멸시효의 진행에 영향을 미치지 않는다.
  \item 수분양자의 건물에 관한 소유권이전등기청구권에 있어서 그 목적물인 건물이 완공되지 아니하여 이를 행사할 수 없었다는 사유는 법률상의 장애사유에 해당하지 않는다.
  \item 정지조건부 권리의 경우에는 조건 미성취의 동안은 권리를 행사할 수 없으므로 소멸시효가 진행되지 않는다.
  \item 소유권이전등기 말소등기의무의 이행불능으로 인한 전보배상청구권의 소멸시효는 말소등기의무가 이행불능 상태에 빠진 때부터 진행된다.
  \item 본래의 소멸시효 기산일과 당사자가 주장하는 기산일이 서로 다른 경우에는 원칙적으로 법원은 당사자가 주장하는 기산일을 기준으로 소멸시효를 계산하여야 한다.
\end{choices}
\end{question}
\answer{2}
